\subsection{A property special to the essential upper integral}

The following result will be used frequently in the sequel. In the statement, one cannot replace essential upper integrals by ordinary upper integrals (see Exer. 4).

PROPOSITION 11. --- Let \((\lambda_{\alpha})_{\alpha\in\mathrm{A}}\) be a family of positive measures on T, directed for the relation \(\leqslant\) and having a supremum \(\lambda\) in \(\mathcal{M}(\mathrm{T})\) . Then, for every numerical function \(f\geqslant0\) ,

\[
\lambda^ {\bullet} (f) = \sup _ {\alpha \in A} \lambda_ {\alpha} ^ {\bullet} (f).\tag{5}
\]

When \(f\) belongs to \(\mathcal{K}(\mathrm{T})\) , this relation reduces to the definition of the supremum of a directed set in \(\mathcal{M}(\mathrm{T})\) (Ch. II, §2, No.~2, Lemma). Suppose next that \(f \leqslant g\) for some function \(g \in \mathcal{K}_+\) (in other words, that \(f\) is bounded and is zero outside a compact set \(\mathbf{K}\) ); let \(\alpha\) be an index such that \(\lambda_{\alpha}(g) \geqslant \lambda(g) - \varepsilon\) , where \(\varepsilon\) is a number \(>0\) ; since the measure \(\nu = \lambda - \lambda_{\alpha}\) is positive, we have \(\nu^{*}(f) \leqslant \nu(g) \leqslant \varepsilon\) , or \(\lambda_{\alpha}^{*}(f) \geqslant \lambda^{*}(f) - \varepsilon\) (Ch. IV, §1, No.~3, Prop. 15). It follows (since \(\varepsilon\) is arbitrary) that the second member of (5) is \(\geqslant\) the first; the reverse inequality being obvious, (5) is established for the special case under consideration. Next, suppose that \(f\) is zero outside \(\mathbf{K}\) but is not necessarily bounded, and set \(f_n = \inf(f, n)\) for every integer \(n\) . Then

\[
\lambda^ {\bullet} (f) = \sup _ {n \in \mathbf {N}} \lambda^ {\bullet} (f _ {n}) = \sup _ {n \in \mathbf {N}} \sup _ {\alpha \in \mathrm{A}} \lambda_ {\alpha} ^ {\bullet} (f _ {n}) = \sup _ {\alpha \in \mathrm{A}} \sup _ {n \in \mathbf {N}} \lambda_ {\alpha} ^ {\bullet} (f _ {n}) = \sup _ {\alpha \in \mathrm{A}} \lambda_ {\alpha} ^ {\bullet} (f).
\]

Finally, with no restriction made on \(f\) , denoting by \(\mathfrak{K}\) the set of compact subsets of \(T\) we have

\[
\begin{array}{c} \lambda^ {\bullet} (f) = \sup _ {K \in \mathfrak {K}} \lambda^ {\bullet} (f \varphi_ {K}) = \sup _ {K \in \mathfrak {K}} \sup _ {\alpha \in A} \lambda_ {\alpha} ^ {\bullet} (f \varphi_ {K}) \\ = \sup _ {\alpha \in A} \sup _ {K \in \mathfrak {K}} \lambda_ {\alpha} ^ {\bullet} (f \varphi_ {K}) = \sup _ {\alpha \in A} \lambda_ {\alpha} ^ {\bullet} (f). \end{array}
\]

COROLLARY 1. --- For a subset N of T to be locally \(\lambda\) -negligible, it is necessary and sufficient that N be locally \(\lambda_{\alpha}\) -negligible for every \(\alpha \in A\) .

COROLLARY 2. --- For a mapping g of T into a topological space G to be \(\lambda\) -measurable, it is necessary and sufficient that it be \(\lambda_{\alpha}\) -measurable for every \(\alpha \in A\) .

The condition is obviously necessary, since \(\lambda_{\alpha} \leqslant \lambda\) for every \(\alpha\) (Ch. IV, §1, No.~3, Prop. 15). Conversely, suppose that \(g\) is \(\lambda_{\alpha}\) -measurable for all \(\alpha\) , denote by \(\mathfrak{K}\) the set of compact subsets K of T such that \(g|_{\mathbf{K}}\) is continuous, and let L be a compact set such that \(\mathrm{L} \cap \mathrm{K}\) is \(\lambda\) -negligible for every \(\mathrm{K} \in \mathfrak{K}\) . Since the set \(\mathfrak{K}\) is \(\lambda_{\alpha}\) -dense, L is \(\lambda_{\alpha}\) -negligible for every \(\alpha\) (Ch. IV, §5, No.~8, Prop. 12), hence is \(\lambda\) -negligible (Cor. 1). It follows that \(\mathfrak{K}\) is \(\lambda\) -dense and that \(g\) is \(\lambda\) -measurable (Ch. IV, §5, No.~10, Prop. 15).

